\documentclass[../main.tex]{subfiles}
\begin{document}
\chapter*{TÓM TẮT NỘI DUNG ĐỒ ÁN}
\thispagestyle{plain}

Trong bối cảnh hầu hết sinh viên hiện nay đều đã sử dụng Zalo, việc triển khai ứng dụng điểm
danh trực tiếp trên nền tảng Zalo Mini App mang lại sự thuận tiện lớn: người dùng có thể điểm
danh ngay mà không cần cài đặt thêm bất kỳ ứng dụng mới nào. Trong khi đó, các phương pháp
điểm danh truyền thống như gọi tên hay ký danh sách giấy lại tiêu tốn nhiều thời gian đối với
các lớp học đông và khó bảo đảm sinh viên thực sự có mặt.

Đồ án trình bày việc xây dựng hệ thống điểm danh thông minh mang tên \textit{Zimo Check-in}
trên nền tảng Zalo Mini App. Quy trình điểm danh được thiết kế gồm bốn bước: quét mã QR động
được ký bằng HMAC-SHA256 và làm mới sau mỗi 30 giây, xác minh khuôn mặt ngay trên thiết bị,
trao đổi mã xác nhận ngang hàng giữa các sinh viên, và tổng hợp lại thành điểm tin cậy. Việc
kết hợp nhiều lớp xác minh bổ trợ này giúp bảo đảm tính chính xác và trung thực của kết quả
điểm danh mà không đòi hỏi phần cứng chuyên dụng đắt tiền.

Về mặt hệ thống, sản phẩm gồm một Mini App dành cho sinh viên và giảng viên, một trang quản
trị web phục vụ phòng đào tạo, và các hàm xử lý phía máy chủ trên nền tảng Firebase. Hệ thống
còn cung cấp nhiều tính năng hỗ trợ như cổng máy chiếu hiển thị mã QR cỡ lớn cho lớp đông,
giám sát điểm danh theo thời gian thực và quản lý đơn xin nghỉ. Kết quả thử nghiệm cho thấy hệ
thống vận hành ổn định và đáp ứng tốt các yêu cầu về chức năng cũng như bảo mật đã đặt ra.

\vspace{1em}
\noindent\textbf{Từ khoá:} \textit{điểm danh thông minh, Zalo Mini App, QR động, HMAC-SHA256, xác minh khuôn mặt, Firebase.}

\vspace{1em}
\begin{flushright}
Sinh viên thực hiện

\vspace{2em}
\textit{(Ký và ghi rõ họ tên)}
\end{flushright}
\end{document}
